\SourceLeaf{65}%
\begin{center}
\textit{THEOREMATIS FUNDAMENTALIS IN THEORIA RESIDUORUM QUADRATICORUM}\par
\textit{DEMONSTRATIO SEXTA.}
\end{center}

\Article{1}

\textsc{Theorema.}\quad \textit{Designante $p$ numerum primum (positivum imparem), $n$ integrum} \textit{positivum per $p$ non divisibilem, $x$ quantitatem indeterminatam, functio}

\[
\begin{aligned}
1+x^n+x^{2n}+x^{3n}+\text{etc.}+x^{np-n}
\end{aligned}
\]

\textit{divisibilis erit per}

\[
\begin{aligned}
1+x+xx+x^3+\text{etc.}+x^{p-1}
\end{aligned}
\]

\textsc{Demonstr.}\quad Accipiatur integer positivus $g$ ita ut fiat $gn\equiv1\quad(\text{mod. }p)$, statuaturque $gn=1+hp$.\quad Tunc erit

{\small
\[
\begin{aligned}
\frac{1+x^n+x^{2n}+x^{3n}+\text{etc.}\,x^{np-n}}{1+x+xx+x^3+\text{etc.}\,x^{p-1}} &=\frac{(1-x^{np})(1-x)}{(1-x^n)(1-x^p)}\\&=\frac{(1-x^{np})(1-x^{gn}-x+x^{hp+1})}{(1-x^n)(1-x^p)}\\
&=\frac{1-x^{np}}{1-x^p}\cdot\frac{1-x^{gn}}{1-x^n}-\frac{x(1-x^{np})}{1-x^n}\cdot\frac{1-x^{hp}}{1-x^p}
\end{aligned}
\]
}

adeoque manifesto functio integra.\quad Q. E. D.

Quaelibet itaque functio integra ipsius $x$ per $\frac{1-x^{np}}{1-x^n}$ divisibilis. etiam divisibilis erit per $\frac{1-x^p}{1-x}$.

\Article{2}

Designet $\alpha$ radicem primitivam positivam pro modulo $p$, i. e. sit $\alpha$ integer positivus talis, ut residua minima positiva potestatum $1,\alpha,\alpha\alpha,\alpha^3\dotsfive\alpha^{p-2}$ secundum modulum $p$ sine respectu ordinis cum numeris $1,2,3,4\dotsfive p-1$ identica fiant.\quad Designando porro per $fx$ functionem

\[
\begin{aligned}
x+x^\alpha+x^{\alpha\alpha}+x^{\alpha^3}+\text{etc.}+x^{\alpha^{p-2}}+1
\end{aligned}
\]

patet, $fx-1- x-xx-x^3-\text{etc.}-x^{p-1}$ divisibilem fore per $1-x^p$, adeoque a potiori per $\frac{1-x^p}{1-x}=1+x+xx+x^3+\text{etc.}+x^{p-1}$, per quam itaque functionem ipsa quoque $fx$ divisibilis erit.\quad Hinc vero sequitur, quum $x$ exprimat quantitatem indeterminatam, esse quoque $f(x^n)$ divisibilem per $\frac{1-x^{np}}{1-x^n}$, et proin (art. praec.) etiam per $\frac{1-x^p}{1-x}$, quoties quidem $n$ sit integer per $p$ non divisibilis. Contra, quoties $n$ est integer per $p$ divisibilis, singulae partes functionis $f(x^n)$ uni%
